% Editable transcription: official PDF pp.1:8--1:10 and pp.1:14--1:16.
\paragraph{Definition 2.2 (Max-divergence).} Let $\mu$ and $\widetilde\mu$ be discrete probability distributions over some set $\Omega$. The max-divergence of $\widetilde\mu$ from $\mu$ is defined by
\[D_\infty(\widetilde\mu\Vert\mu)=\ln\left(\max_{x\in\Omega}\left(\frac{\Pr[\widetilde{\mathbf x}=x]}{\Pr[\mathbf x=x]}\right)\right),\]
where $\mathbf x\sim\mu$ and $\widetilde{\mathbf x}\sim\widetilde\mu$.
Max-divergence and total variation distance are related by the following lemma.
\paragraph{Lemma 2.3 (Low max-divergence $\Rightarrow$ low total variation distance).} Let $\mu$ and $\widetilde\mu$ be discrete probability distributions over the same set $\Omega$. Let $\varepsilon=D_\infty(\widetilde\mu\Vert\mu)$. There exists a probability distribution $\mu'$ such that $\mu$ can be written as a convex combination $\mu=(1-\varepsilon)\cdot\widetilde\mu+\varepsilon\cdot\mu'$. Moreover, the total variation distance between $\mu$ and $\widetilde\mu$ is at most $\varepsilon$.
\begin{proof}
If $\varepsilon=0$, the lemma is trivial, so assume $\varepsilon>0$. For each $x\in\Omega$, define
\[p(x)=\frac{\Pr[\mathbf x=x]-(1-\varepsilon)\Pr[\widetilde{\mathbf x}=x]}{\varepsilon},\]
where $\mathbf x\sim\mu$ and $\widetilde{\mathbf x}\sim\widetilde\mu$. Then $\sum_{x\in\Omega}p(x)=1$. Furthermore, $p(x)\ge0$, because
\[(1-\varepsilon)\Pr[\widetilde{\mathbf x}=x]\le(1-\varepsilon)\cdot e^\varepsilon\cdot\Pr[\mathbf x=x]\le\Pr[\mathbf x=x].\]
Therefore, $p(\cdot)$ is a probability mass function, and we can let $\mu'$ be the corresponding probability distribution. For the ``moreover'' part, observe that for any $S\subseteq\Omega$, we have
\[\Pr[\mathbf x\in S]=(1-\varepsilon)\cdot\Pr[\widetilde{\mathbf x}\in S]+\varepsilon\cdot\Pr[\mathbf x'\in S],\]
where $\mathbf x'\sim\mu'$. Therefore,
\[\Pr[\mathbf x\in S]\le\Pr[\widetilde{\mathbf x}\in S]+\varepsilon\cdot\Pr[\mathbf x'\in S]\le\Pr[\widetilde{\mathbf x}\in S]+\varepsilon,\]
and
\[\Pr[\mathbf x\in S]\ge(1-\varepsilon)\cdot\Pr[\widetilde{\mathbf x}\in S]\ge\Pr[\widetilde{\mathbf x}\in S]-\varepsilon.\]
\end{proof}
\paragraph{Definition 2.6 (Generalized restriction).} A generalized restriction is a function $\pi:\{\pm1\}^r\to\{\pm1\}^n$. If $f:\{\pm1\}^n\to\{\pm1\}$ is a Boolean function, then we define $g|_\pi$ to be the composition $g\circ\pi$. That is, $g|_\pi:\{\pm1\}^r\to\{\pm1\}$ is given by $g|_\pi(x)=g(\pi(x))$.
Traditional restrictions can be viewed as a special case of generalized restrictions as follows.
\paragraph{Definition 2.7 (Traditional restrictions as generalized restrictions).} A restriction is a string $\rho\in\{+1,-1,\star\}^n$. We identify $\rho$ with a generalized restriction $\pi:\{\pm1\}^r\to\{\pm1\}^n$, where $r=|\rho^{-1}(\star)|$, as follows. Given $y\in\{\pm1\}^r$, we let $\pi(y)$ be $\rho$, except that the $i$-th star is replaced with $y_i$ for every $i\in[r]$.
Next, we consider distributions over generalized restrictions, and we explain how to interpret the tensor product of such distributions.
\paragraph{Definition 2.8 (Tensor product of generalized restriction distributions).} Let $r,n\in\mathbb N$, and let $\mathcal R$ be a distribution over generalized restrictions $\pi:\{\pm1\}^r\to\{\pm1\}^n$. Let $\boldsymbol\pi_1,\ldots,\boldsymbol\pi_t$ be independent samples from $\mathcal R$, and define $\vec{\boldsymbol\pi}:\{\pm1\}^{rt}\to\{\pm1\}^{nt}$ by concatenating, i.e.,
\[\vec{\boldsymbol\pi}(y^{(1)},\ldots,y^{(t)})=(\boldsymbol\pi_1(y^{(1)}),\ldots,\boldsymbol\pi_t(y^{(t)})).\]
Then the tensor product $\mathcal R^{\otimes t}$ is the distribution of the random variable $\vec{\boldsymbol\pi}$.
% STATEMENT_BEGIN
\paragraph{Lemma 4.2 (Tighter XOR lemma for the random simplification method).}
Let $n,t\in\mathbb N$ and let $h:\{\pm1\}^n\to\{\pm1\}$ be a Boolean function. Let $\mathcal C$ be a class of Boolean functions $g:\{\pm1\}^{nt}\to\{\pm1\}$ that is closed under restrictions.\footnote{Every function in $\mathcal C$ has domain $\{\pm1\}^{nt}$. When we say that $\mathcal C$ is closed under restrictions, we are thinking of a restriction of $g\in\mathcal C$ as another function on $nt$ bits that ignores some of its input variables.} Let $r\in\mathbb N$, let $\mathcal R$ be a distribution over generalized restrictions $\pi:\{\pm1\}^r\to\{\pm1\}^n$, and let $\mu$ be a distribution over $\{\pm1\}^r$. Let $\varepsilon>0$, and assume the following.
\begin{enumerate}
\item (The distribution $\mu$ approximately completes $\mathcal R$ to the uniform distribution.) If we sample $\boldsymbol\pi\sim\mathcal R$ and $\mathbf y\sim\mu$ independently, and we sample $\mathbf x\in\{\pm1\}^n$ uniformly at random, then $D_\infty(\boldsymbol\pi(\mathbf y)\Vert\mathbf x)\le\varepsilon$.
\item (The class $\mathcal C$ simplifies under $\mathcal R^{\otimes t}$.) For every $g\in\mathcal C$ and every $D\in\mathbb N$, we have
\[\Pr_{\vec{\boldsymbol\pi}\sim\mathcal R^{\otimes t}}[\DT(g|_{\vec{\boldsymbol\pi}})\ge D]\le2^{t-D}.\]
\item (The function $h$ retains structure under $\mathcal R$.) For every $D\in\mathbb N$ and every $\pi\in\operatorname{Supp}(\mathcal R)$, we have
\[\Corr_\mu(h|_\pi,\DT[D])\le\varepsilon\cdot2^{D/3}.\]
\end{enumerate}
Then $\Corr(\mathcal C,h^{\oplus t})\le O(\varepsilon)^t$.
% STATEMENT_END
% PROOF_BEGIN
\begin{proof}
Fix any $g\in\mathcal C$. Our job is to analyze the correlation between $g$ and $h^{\oplus t}$ under a uniform random input. By Lemma 2.3, we can sample a uniform random input by the following procedure.
\begin{enumerate}
\item Sample $\vec{\boldsymbol\pi}=(\boldsymbol\pi_1,\ldots,\boldsymbol\pi_t)\sim\mathcal R^{\otimes t}$.
\item Sample $\vec{\mathbf y}=(\mathbf y^{(1)},\ldots,\mathbf y^{(t)})\sim\mu^{\otimes t}$.
\item Sample $\vec{\mathbf e}=(\mathbf e^{(1)},\ldots,\mathbf e^{(t)})\sim(\mu')^{\otimes t}$, where $\mu'$ is the distribution over $\{\pm1\}^n$ from Lemma 2.3.
\item Sample $\mathbf I\subseteq[t]$ where $\Pr[i\in\mathbf I]=1-\varepsilon$ independently for every $i$.
\item Output the string $\vec{\mathbf x}=(\mathbf x^{(1)},\ldots,\mathbf x^{(t)})\in\{\pm1\}^{nt}$, where
\[\mathbf x^{(i)}=\begin{cases}\boldsymbol\pi_i(\mathbf y^{(i)})&\text{if }i\in\mathbf I\\\mathbf e^{(i)}&\text{if }i\notin\mathbf I.\end{cases}\]
\end{enumerate}
Let $\mathbf g:\{\pm1\}^{nt}\to\{\pm1\}$ be the function obtained from $g$ by plugging $\mathbf e^{(i)}$ into each block $i\notin\mathbf I$ and leaving the blocks in $\mathbf I$ alive. Since $\mathbf g$ ignores the variables in blocks outside $\mathbf I$, we have
\[g(\vec{\mathbf x})=\mathbf g|_{\vec{\boldsymbol\pi}}(\vec{\mathbf y}).\]
Similarly, define $\mathbf h:\{\pm1\}^{nt}\to\{\pm1\}$ by the formula
\[\mathbf h(\vec x)=\left(\prod_{i\in\mathbf I}h_i(x^{(i)})\right)\cdot\left(\prod_{i\notin\mathbf I}h_i(\mathbf e^{(i)})\right),\]
so that $h^{\oplus t}(\vec{\mathbf x})=\mathbf h|_{\vec{\boldsymbol\pi}}(\vec{\mathbf y})$. That way,
\[\Corr(g,h^{\oplus t})=\mathbb E[g(\vec{\mathbf x})\cdot h^{\oplus t}(\vec{\mathbf x})]=\mathbb E_{\mathbf I,\vec{\mathbf e},\vec{\boldsymbol\pi}}[\Corr_{\mu^{\otimes t}}(\mathbf g|_{\vec{\boldsymbol\pi}},\mathbf h|_{\vec{\boldsymbol\pi}})].\]
Let $\mathbf D=\lfloor\DT(\mathbf g|_{\vec{\boldsymbol\pi}})/|\mathbf I|\rfloor$. Then
\[\mathbb E_{\mathbf I,\vec{\mathbf e},\vec{\boldsymbol\pi}}[\Corr_{\mu^{\otimes t}}(\mathbf g|_{\vec{\boldsymbol\pi}},\mathbf h|_{\vec{\boldsymbol\pi}})]\le\mathbb E_{\mathbf I,\vec{\mathbf e},\vec{\boldsymbol\pi}}[\Corr_{\mu^{\otimes t}}(\mathbf h|_{\vec{\boldsymbol\pi}},\DT[(\mathbf D+1)\cdot|\mathbf I|])].\]
Let $\boldsymbol\pi_{\mathbf I}=(\boldsymbol\pi_i)_{i\in\mathbf I}$, and define $h_I:\{\pm1\}^{n|I|}\to\{\pm1\}$ by $h_I((x^{(i)})_{i\in I})=\prod_{i\in I}h(x^{(i)})$. Then for any fixing of $\mathbf I,\vec{\mathbf e},\vec{\boldsymbol\pi}$, we have
\[\Corr_{\mu^{\otimes t}}(\mathbf h|_{\vec{\boldsymbol\pi}},\DT[(\mathbf D+1)\cdot|\mathbf I|])=\Corr_{\mu^{\otimes|\mathbf I|}}(h_{\mathbf I}|_{\boldsymbol\pi_{\mathbf I}},\DT[(\mathbf D+1)\cdot|\mathbf I|]).\]
Now we apply Lemma 3.2. For each $i\in\mathbf I$ and each $D\in\mathbb N$, we have
\[\Corr_\mu(h|_{\boldsymbol\pi_i},\DT[D])\le\varepsilon\cdot2^{D/3}.\]
Furthermore, the function $\varepsilon(D)=\varepsilon\cdot2^{D/3}$ is log-concave. Therefore, Lemma 3.2 guarantees that
\[\Corr_{\mu^{\otimes|\mathbf I|}}(h_{\mathbf I}|_{\boldsymbol\pi_{\mathbf I}},\DT[(\mathbf D+1)\cdot|\mathbf I|])\le O(\varepsilon\cdot2^{2\cdot(\mathbf D+1)/3})^{|\mathbf I|}=O(\varepsilon\cdot2^{2\mathbf D/3})^{|\mathbf I|}.\]
Thus, overall, we get
\[\Corr(g,h^{\oplus t})\le\mathbb E_{\mathbf I,\vec{\mathbf e},\vec{\boldsymbol\pi}}\left[O\left(\varepsilon\cdot2^{2\mathbf D/3}\right)^{|\mathbf I|}\right].\]
Now consider any fixing of $\mathbf I$ and $\vec{\mathbf e}$. Since $\mathcal C$ is closed under restrictions, $\mathbf g\in\mathcal C$. Therefore, our simplification assumption tells us that for every $D\in\mathbb N$, we have
\[\Pr_{\vec{\boldsymbol\pi}}[\mathbf D=D]\le2^{t-D\cdot|\mathbf I|}.\]
Consequently,
\begin{align*}
\mathbb E_{\vec{\boldsymbol\pi}}\left[O\left(\varepsilon\cdot2^{2\mathbf D/3}\right)^{|\mathbf I|}\right]&=\sum_{D=0}^\infty\Pr_{\vec{\boldsymbol\pi}}[\mathbf D=D]\cdot O\left(\varepsilon\cdot2^{2D/3}\right)^{|\mathbf I|}\\
&\le2^t\cdot\sum_{D=0}^\infty O\left(\varepsilon\cdot2^{-D/3}\right)^{|\mathbf I|}\\
&\le2^t\cdot O\left(\sum_{D=0}^\infty\varepsilon\cdot2^{-D/3}\right)^{|\mathbf I|}\\
&=2^t\cdot O(\varepsilon)^{|\mathbf I|}.
\end{align*}
Therefore, our overall bound is given by
\begin{align*}
\Corr(g,h^{\oplus t})&\le\mathbb E_{\mathbf I}[2^t\cdot O(\varepsilon)^{|\mathbf I|}]=2^t\cdot\sum_{I\subseteq[t]}\Pr[\mathbf I=I]\cdot O(\varepsilon)^{|I|}\\
&=2^t\cdot\sum_{I\subseteq[t]}(1-\varepsilon)^{|I|}\cdot\varepsilon^{t-|I|}\cdot O(\varepsilon)^{|I|}\\
&\le O(\varepsilon)^t.
\end{align*}
\end{proof}
% PROOF_END
